
\setcounter{subsection}{3}
\setcounter{subsubsection}{0}
\begin{sublemme}
\label{XIII.4.3.1}
Sean $f\colon X\to S$ un morfismo $0$-acíclico y localmente
$0$-acíclico, y~$g$ una sección de~$f$. Sea~$C$ un grupo constante
finito tal que $(C_X,f)$ sea cohomológicamente propio en dimensión
$\leqslant0$ y que la imagen directa por~$f$ del stack de los torsores bajo
$C_X$ sea constructible. Entonces, para toda sección~$z$ de
$\H^0(S,\R^1f_*C_X)$, puede encontrarse un recubrimiento étale
$S_1$ de~$S$ y,
\ifthenelse{\boolean{orig}}
{si $X_1=X\times_SS_1$, un elemento de}
{si $X_1=X\times_SS_1$ y $f_1\colon X_1\to S_1$ es la proyección, un elemento de}
$\H^1(X_1,C_{X_1})$ cuya imagen por el morfismo canónico
$$
\H^1(X_1,C_{X_1})\to
\ifthenelse{\boolean{orig}}
{\H^0(S_1,\R^1f_*C_{X_1})}
{\H^0(S_1,\R^1f_{1*}C_{X_1})}
$$
sea igual a la restricción de~$z$ a
\ifthenelse{\boolean{orig}}
{$\H^0(S_1,\R^1f_*C_{X_1})$}
{$\H^0(S_1,\R^1f_{1*}C_{X_1})$}.
\end{sublemme}

Para todo esquema~$S'$ étale sobre~$S$, se pone $X'=X\times_SS'$ y
\ifthenelse{\boolean{orig}}
{se denota por~$g'$ (\resp $F'$, etc.) la imagen inversa de~$g$
(\resp $F$, etc.) por el morfismo}
{se denota por~$g'$ la imagen inversa de~$g$ por el morfismo}
\marginpar{419}
$S'\to S$. El prehaz~$G$ sobre~$S$ definido por
\newlength{\tmpRacinet} \setlength{\tmpRacinet}{\textwidth}
\addtolength{\tmpRacinet}{-7em}
$$
G(S') = \left\{\text{\parbox{\tmpRacinet}{clases, salvo isomorfismo, de
		torsores $P$ sobre~$X'$ bajo~$C_{X'}$, provistos de un
		isomorfismo $g'^*P\isomto C_{S'}$ }} \right\}
$$
es entonces un haz. Esto resulta, en efecto, por descenso, del hecho
de que un automorfismo del torsor~$P$ sobre~$X'$ queda bien
determinado por su restricción a~$g'(S')$. Además, se tiene un
morfismo sobreyectivo
$$
\ifthenelse{\boolean{orig}}
{G\to \R^1f_*F.}
{G\to \R^1f_*C_X.}
$$
Sean~$z$ un elemento de~$\H^0(S,\R^1f_*C_X)$ y~$H$ el
subhaz de~$G$ que es la imagen inversa de~$z$. Basta demostrar
que~$H$ es un haz localmente constante y constructible. Como esta
propiedad es local sobre~$S$, puede suponerse que~$z$
proviene de un elemento de~$\H^1(X,C_X)$ representado por
un torsor~$P$ tal que $g^*P$ sea isomorfo a~$C_S$. Dar un
isomorfismo $i\colon g^*P\isomto C_S$ equivale a dar una
sección global de~$\SheafAut_{C_S}(g^*P)$, y dos isomorfismos~$i$
e~$i'$ definen el mismo elemento de
\ifthenelse{\boolean{orig}}
{$G(X)$}
{$G(S)$}
si y solo si
\ifthenelse{\boolean{orig}}
{$ii'^{-1}$}
{$i^{-1}i'$}
es la imagen de un elemento de
$\Aut_{C_X}(P)$. Si se considera la inyección canónica
$$
f_*\SheafAut_{C_X}(P)\to\SheafAut_{C_S}(g^*P)\simeq C,
$$
$H$ se identifica, por tanto, con el cociente de
$\SheafAut_{C_S}(g^*P)$ por~$f_*\SheafAut_{C_X}(P)$. Según
\Ref{XIII.1.16}, $f_*\SheafAut_{C_X}(P)$ es localmente constante;
lo mismo ocurre, pues, con~$H$, lo que concluye la demostración.
